% Part: first-order-logic
% Chapter: completeness
% Section: complete-sets

% Definition of complete consistent sets. Properties of complete sets required
% for completeness proved are in provability.tex in the
% chapter on the proof system used.

\documentclass[../../../include/open-logic-section]{subfiles}

\begin{document}

\iftag{FOL}
      {\olfileid[pt]{fol}{com}{ccs}}
      {\olfileid[pt]{pl}{com}{ccs}}
      
\olsection{Conjuntos Consistentes Completos de \usetoken{P}{sentence}}

\begin{defn}[Conjunto completo]
\ollabel{def:complete-set} Um conjunto~$\Gamma$ de !!{sentence}s é
\emph{!!{complete}} se e somente se, para qualquer !!{sentence}~$!A$, ou $!A \in
\Gamma$ ou $\lnot !A \in \Gamma$.
\end{defn}

\begin{explain}
Conjuntos !!{complete}s de sentenças não deixam questões sem resposta. Para
qualquer !!{sentence}~$!A$, $\Gamma$ ``diz'' se $!A$ é verdadeira ou falsa.  A
importância dos conjuntos !!{complete}s estende-se para além da prova do
teorema da completude. Uma teoria que é completa e
axiomatizável, por exemplo, é sempre decidível.
\end{explain}

\begin{explain}
Conjuntos consistentes !!{complete}s são importantes na prova de completude
já que podemos garantir que todo conjunto consistente de
!!{sentence}s~$\Gamma$ está contido em um conjunto consistente !!{complete}
~$\Gamma^*$.  Um conjunto consistente !!{complete} contém, para cada
!!{sentence}~$!A$, ou $!A$ ou sua negação $\lnot !A$, mas não
ambas. Isso é verdade em particular para \iftag{FOL}{!!{sentence}s
atômicas}{variáveis proposicionais}, de modo que, a partir de um conjunto consistente
!!{complete}\iftag{FOL}{ em uma linguagem adequadamente expandida por
!!{constant}s}{}, podemos construir
\iftag{FOL}{!!a{structure}}{!!a{valuation}} onde a
\iftag{FOL}{interpretação dos !!{predicate}s}{valor de verdade atribuído às
variáveis proposicionais} é definida de acordo com quais
\iftag{FOL}{!!{sentence}s atômicas}{!!{propositional variable}s} estão
em~$\Gamma^*$. Pode-se então mostrar que esta
\iftag{FOL}{!!{structure}}{!!{valuation}} torna verdadeiras todas as
!!{sentence}s em~$\Gamma^*$ (e, portanto, também todas
aquelas em~$\Gamma$). A prova deste último fato requer que
$\lnot !A \in \Gamma^*$ se e somente se $!A \notin \Gamma^*$, $(!A \lor !B) \in
\Gamma^*$ se e somente se $!A \in \Gamma^*$ ou $!B \in \Gamma^*$, etc.
\end{explain}

No que se segue, usaremos com frequência, de modo tácito, as propriedades de
reflexividade, monotonicidade e transitividade de $\Proves$ (ver
\iftag{FOL}{%
  \tagrefs{prfSC/{fol:seq:ptn:sec},prfND/{fol:ntd:ptn:sec},prfAX/{fol:axd:ptn:sec},prfTab/{fol:tab:ptn:sec}}}{%
  \tagrefs{prfSC/{pl:seq:ptn:sec},prfND/{pl:ntd:ptn:sec},prfAX/{pl:axd:ptn:sec},prfTab/{pl:tab:ptn:sec}}}).

\begin{prop}
\ollabel{prop:ccs}
Suponha que $\Gamma$ seja !!{complete} e consistente. Então:
\begin{enumerate}
\item \ollabel{prop:ccs-prov-in} Se $\Gamma \Proves !A$, então $!A \in
  \Gamma$.

\tagitem{prvAnd}{\ollabel{prop:ccs-and} $!A \land !B \in \Gamma$
  se e somente se $!A \in \Gamma$ e $!B \in \Gamma$.}{}

\tagitem{prvOr}{\ollabel{prop:ccs-or} $!A \lor !B \in \Gamma$ se e somente se
  $!A \in \Gamma$ ou $!B \in \Gamma$.}{}

\tagitem{prvIf}{\ollabel{prop:ccs-if} $!A \lif !B \in \Gamma$ se e somente se
  $!A \notin \Gamma$ ou $!B \in \Gamma$.}{}
\end{enumerate}
\end{prop}

\begin{proof}
Suponhamos, para tudo o que se segue, que $\Gamma$ seja !!{complete} e
consistente.
\begin{enumerate}
\item Se $\Gamma \Proves !A$, então $!A \in \Gamma$.

Suponha que $\Gamma \Proves !A$. Suponha, ao contrário, que $!A
\notin \Gamma$. Como $\Gamma$ é !!{complete}, $\lnot !A \in \Gamma$.
Por
\iftag{FOL}{%
  \tagrefs{prfSC/{fol:seq:prv:prop:explicit-inc},
    prfND/{fol:ntd:prv:prop:explicit-inc},
    prfAX/{fol:axd:prv:prop:explicit-inc},
    prfTab/{fol:tab:prv:prop:explicit-inc}}}{%
  \tagrefs{prfSC/{pl:seq:prv:prop:explicit-inc},
    prfND/{pl:ntd:prv:prop:explicit-inc},
    prfAX/{pl:axd:prv:prop:explicit-inc},
    prfTab/{pl:tab:prv:prop:explicit-inc}}},
$\Gamma$ é inconsistente. Isso contradiz a suposição de que
$\Gamma$ é consistente. Logo, não pode ser o caso que $!A \notin
\Gamma$, de modo que $!A \in \Gamma$.

\tagitem{defAnd}{}{%
\iftag{probAnd}{Exercício.}{%
$!A \land !B \in \Gamma$ se e somente se $!A \in \Gamma$ e $!B \in \Gamma$:

Para a direção direta, suponha $!A \land !B \in \Gamma$. Então
por
\iftag{FOL}{%
  \tagrefs{prfSC/{fol:seq:ppr:prop:provability-land},
    prfND/{fol:ntd:ppr:prop:provability-land},
    prfAX/{fol:axd:ppr:prop:provability-land},
    prfTab/{fol:tab:ppr:prop:provability-land}}}{%
  \tagrefs{prfSC/{pl:seq:ppr:prop:provability-land},
    prfND/{pl:ntd:ppr:prop:provability-land},
    prfAX/{pl:axd:ppr:prop:provability-land},
    prfTab/{pl:tab:ppr:prop:provability-land}}}, item~(1),
$\Gamma \Proves !A$ e $\Gamma \Proves !B$. Por
\olref{prop:ccs-prov-in}, $!A \in \Gamma$ e $!B \in \Gamma$, como
requerido.

Para a direção reversa, sejam $!A \in \Gamma$ e $!B \in
\Gamma$. Por
\iftag{FOL}{%
  \tagrefs{prfSC/{fol:seq:ppr:prop:provability-land},
    prfND/{fol:ntd:ppr:prop:provability-land},
    prfAX/{fol:axd:ppr:prop:provability-land},
    prfTab/{fol:tab:ppr:prop:provability-land}}}{%
  \tagrefs{prfSC/{pl:seq:ppr:prop:provability-land},
    prfND/{pl:ntd:ppr:prop:provability-land},
    prfAX/{pl:axd:ppr:prop:provability-land},
    prfTab/{pl:tab:ppr:prop:provability-land}}}, item~(2),
$\Gamma \Proves !A \land !B$. Por \olref{prop:ccs-prov-in}, $!A \land
!B \in \Gamma$.}}

\tagitem{defOr}{}{%
\iftag{probOr}{Exercício.}{%
Primeiro mostramos que, se $!A \lor !B \in \Gamma$, então $!A \in
\Gamma$ ou $!B \in \Gamma$. Suponha $!A \lor !B \in \Gamma$ mas $!A
\notin \Gamma$ e $!B \notin \Gamma$.  Como $\Gamma$ é
!!{complete}, $\lnot !A \in \Gamma$ e $\lnot !B \in \Gamma$. Por
\iftag{FOL}{%
  \tagrefs{prfSC/{fol:seq:ppr:prop:provability-lor},
    prfND/{fol:ntd:ppr:prop:provability-lor},
    prfAX/{fol:axd:ppr:prop:provability-lor},
    prfTab/{fol:tab:ppr:prop:provability-lor}}}{%
  \tagrefs{prfSC/{pl:seq:ppr:prop:provability-lor},
    prfND/{pl:ntd:ppr:prop:provability-lor},
    prfAX/{pl:axd:ppr:prop:provability-lor},
    prfTab/{pl:tab:ppr:prop:provability-lor}}},
item (1), $\Gamma$ é inconsistente, uma contradição. Logo, $!A
\in \Gamma$ ou $!B \in \Gamma$.

Para a direção reversa, suponha que $!A \in \Gamma$ ou $!B \in
\Gamma$. Por
\iftag{FOL}{%
  \tagrefs{prfSC/{fol:seq:ppr:prop:provability-lor},
    prfND/{fol:ntd:ppr:prop:provability-lor},
    prfAX/{fol:axd:ppr:prop:provability-lor},
    prfTab/{fol:tab:ppr:prop:provability-lor}}}{%
  \tagrefs{prfSC/{pl:seq:ppr:prop:provability-lor},
    prfND/{pl:ntd:ppr:prop:provability-lor},
    prfAX/{pl:axd:ppr:prop:provability-lor},
    prfTab/{pl:tab:ppr:prop:provability-lor}}}, item (2),
$\Gamma \Proves !A \lor !B$. Por \olref{prop:ccs-prov-in}, $!A \lor
!B \in \Gamma$, como requerido.}}

\tagitem{defIf}{}{%
\iftag{probIf}{Exercício.}{%
Para a direção direta, suponha $!A \lif !B \in \Gamma$ e suponha,
ao contrário, que $!A \in \Gamma$ e $!B \notin \Gamma$. Sob essas
suposições, $!A \lif !B \in \Gamma$ e $!A \in \Gamma$. Por
\iftag{FOL}{%
  \tagrefs{prfSC/{fol:seq:ppr:prop:provability-lif},
    prfND/{fol:ntd:ppr:prop:provability-lif},
    prfAX/{fol:axd:ppr:prop:provability-lif},
    prfTab/{fol:tab:ppr:prop:provability-lif}}}{%
  \tagrefs{prfSC/{pl:seq:ppr:prop:provability-lif},
    prfND/{pl:ntd:ppr:prop:provability-lif},
    prfAX/{pl:axd:ppr:prop:provability-lif},
    prfTab/{pl:tab:ppr:prop:provability-lif}}}, item (1),
$\Gamma \Proves !B$. Mas então, por \olref{prop:ccs-prov-in}, $!B \in
\Gamma$, contradizendo a suposição de que $!B \notin \Gamma$.

Para a direção reversa, considere primeiro o caso em que $!A \notin
\Gamma$. Como $\Gamma$ é !!{complete}, $\lnot !A \in \Gamma$.  Por
\iftag{FOL}{%
  \tagrefs{prfSC/{fol:seq:ppr:prop:provability-lif},
    prfND/{fol:ntd:ppr:prop:provability-lif},
    prfAX/{fol:axd:ppr:prop:provability-lif},
    prfTab/{fol:tab:ppr:prop:provability-lif}}}{%
  \tagrefs{prfSC/{pl:seq:ppr:prop:provability-lif},
    prfND/{pl:ntd:ppr:prop:provability-lif},
    prfAX/{pl:axd:ppr:prop:provability-lif},
    prfTab/{pl:tab:ppr:prop:provability-lif}}}, item (2),
$\Gamma \Proves !A \lif !B$. Novamente por \olref{prop:ccs-prov-in}, obtemos
que $!A \lif !B \in \Gamma$, como requerido.

Agora considere o caso em que $!B \in \Gamma$.  Por
\iftag{FOL}{%
  \tagrefs{prfSC/{fol:seq:ppr:prop:provability-lif},
    prfND/{fol:ntd:ppr:prop:provability-lif},
    prfTab/{fol:tab:ppr:prop:provability-lif},
    prfAX/{fol:axd:ppr:prop:provability-lif}}}{%
  \tagrefs{prfSC/{pl:seq:ppr:prop:provability-lif},
    prfND/{pl:ntd:ppr:prop:provability-lif},
    prfTab/{pl:tab:ppr:prop:provability-lif},
    prfAX/{pl:axd:ppr:prop:provability-lif}}},
item (2) novamente, $\Gamma \Proves !A \lif !B$. Por
\olref{prop:ccs-prov-in}, $!A \lif !B \in \Gamma$.}}
\end{enumerate}
\end{proof}

\tagprob{FOL}
\begin{prob}
Complete a prova de \olref[fol][com][ccs]{prop:ccs}.
\end{prob}
\tagendprob

\tagprob{notFOL}
\begin{prob}
Complete a prova de \olref[pl][com][ccs]{prop:ccs}.
\end{prob}
\tagendprob

\end{document}
